

\documentclass{article}

\usepackage{fontsetup}
\usepackage[fontsize=14pt]{fontsize}
\usepackage[ukrainian]{babel}
\usepackage{microtype}
\usepackage{indentfirst}
\usepackage{tabularray}
\usepackage{xcolor}
\usepackage{graphicx}
\graphicspath{{pic/}}
\usepackage{wrapstuff}
\usepackage{amsmath}
\usepackage{physics}
\usepackage{rotating}
\usepackage{multirow}% Гадаю ця бібліотека не є обов'язковою, якщо є tabulararray
\usepackage{float}
\usepackage[style=gost-numeric,
autolang=other]{biblatex}
\addbibresource{bib/bibl.bib}

\usepackage[
colorlinks=true,
urlcolor = blue, %Colour for external hyperlinks
linkcolor  = red, %Colour of internal links
citecolor  = green, %Colour of citations
bookmarks = true,
bookmarksnumbered=true,
unicode,
linktoc = all,
hypertexnames=false,
pdftoolbar=false,
pdfpagelayout=TwoPageRight,
pdfauthor={Sokalskyi B.T},
pdfdisplaydoctitle=true,
pdfencoding=auto
]%
{hyperref}


\usepackage[%
%    showframe,
a4paper,%
footskip=1cm,%
headsep=0.3cm,%
top=2cm, %поле зверху
bottom=2cm, %поле знизу
left=2.5cm, %поле ліворуч
right=1cm, %поле праворуч
]{geometry}

%-----------Підключення необхідних графічних бібліотек---------
\usepackage{tikz}
\usepackage{pgfplots, pgfplotstable}
\pgfplotsset{compat=newest}
\usetikzlibrary{intersections, calc, plotmarks, fpu}
\usepackage{booktabs,colortbl}% Оскільки pgfplotstabe працює з tabular, тому треба використати ці пакети для розфарбовування таблиці
